\DocumentMetadata{
  pdfversion=2.0,
  pdfstandard=ua-2,
  lang=en-US,
  tagging=on,
  tagging-setup={math/setup=mathml-SE,tabsorder=structure}
}
\documentclass[11pt,letterpaper]{article}
\usepackage[margin=0.68in,includefoot]{geometry}
\usepackage{newcomputermodern}
\usepackage{amsmath,amscd,mathtools}
\usepackage{tikz,adjustbox}
\providecommand{\square}{\mdlgwhtsquare}
\usepackage[hidelinks]{hyperref}
\hypersetup{
  pdftitle={Numerical Analysis PhD Exam, September 2011},
  pdfauthor={University of Florida Department of Mathematics},
  pdfsubject={Graduate mathematics examination},
  pdfdisplaydoctitle=true,
  bookmarksopen=true
}
\setcounter{secnumdepth}{0}
\setlength{\parindent}{0pt}
\setlength{\parskip}{0.28em}
\setlength{\emergencystretch}{3em}
\setlength{\leftmargini}{2.15em}
\setlength{\leftmarginii}{2.65em}
\setlength{\leftmarginiii}{3em}
\renewcommand{\labelenumi}{\textbf{\theenumi.}}
\pagestyle{empty}
\raggedbottom
\allowdisplaybreaks
\newenvironment{examproblems}[1]{%
  \begin{enumerate}%
  \setcounter{enumi}{#1}%
  \setlength{\topsep}{0.35em}%
  \setlength{\partopsep}{0pt}%
  \setlength{\itemsep}{0.48em}%
  \setlength{\parsep}{0.18em}%
}{\end{enumerate}}
\newenvironment{examsubparts}{%
  \begin{enumerate}%
  \setlength{\topsep}{0.18em}%
  \setlength{\partopsep}{0pt}%
  \setlength{\itemsep}{0.16em}%
  \setlength{\parsep}{0.1em}%
}{\end{enumerate}}
\newenvironment{examinstructions}{%
  \par\begingroup\itshape
}{%
  \par\endgroup\vspace{0.18em}
}
\makeatletter
\renewcommand\section{\@startsection{section}{1}{0pt}%
  {-1.25ex plus -.25ex minus -.1ex}%
  {0.55ex plus .1ex}%
  {\normalfont\large\bfseries}}
\renewcommand{\@maketitle}{%
  \newpage\null\vskip -1em%
  {\footnotesize\bfseries
    \href{https://www.ufl.edu/}{University of Florida}, \href{https://math.ufl.edu/}{Department of Mathematics}\par}%
  \vskip -0.5em%
  \tagstructbegin{tag=Title}%
    {\LARGE\bfseries\href{https://gma.math.ufl.edu/exams/numerical-analysis-phd/}{Numerical Analysis PhD Exam}, September 2011\par}%
  \tagstructend%
  \vskip 0.25em%
  \tagmcbegin{artifact}%
    \hrule height 1pt%
  \tagmcend%
  \vskip 1em%
}
\makeatother
\title{Numerical Analysis PhD Exam, September 2011}
\author{}
\date{}
\begin{document}
\maketitle
\thispagestyle{empty}
\begin{examinstructions}
Answer any 8 questions.
\end{examinstructions}
\begin{examproblems}{0}
\item
This problem has two parts.
\begin{examsubparts}
\item[\textnormal{(a)}]
For any matrices~\(A\) and~\(B\), show that the nonzero eigenvalues of \(AB\) and \(BA\) are the same.
\item[\textnormal{(b)}]
If \(AB\) is normal, \(\lVert\cdot\rVert_2\) is the 2-norm of a matrix, and \(\lVert\cdot\rVert\) is an induced matrix norm, then show that \(\lVert AB\rVert_2\leq \lVert BA\rVert\).
\end{examsubparts}
\item
This problem has two parts.
\begin{examsubparts}
\item[\textnormal{(a)}]
Suppose~\(p\) and \(q\in\mathbb{R}\) with~\(p\) and~\(q\) positive and \(p^{-1}+q^{-1}=1\). Show that for any matrix \(A\in\mathbb{C}^{m\times n}\), we have \(\lVert A\rVert_p=\lVert A^*\rVert_q\).
\item[\textnormal{(b)}]
Prove that 
\[
\lVert A\rVert_2^2\leq \lVert A\rVert_p\lVert A\rVert_q
\]
 for any \(A\in\mathbb{C}^{m\times n}\) and any positive~\(p\) and \(q\in\mathbb{R}\) with \(p^{-1}+q^{-1}=1\).
\end{examsubparts}
\item
Let~\(P\) and~\(Q\) be two \(m\times m\) orthogonal projectors. Prove that \(\lVert P-Q\rVert_2\leq 1\).
\item
Let~\(P\) and~\(Q\) be Hermitian positive definite matrices. Prove that 
\[
x^*Px\leq x^*Qx\quad\text{for all }x\in\mathbb{C}^n
\]
 if and only if 
\[
x^*Q^{-1}x\leq x^*P^{-1}x\quad\text{for all }x\in\mathbb{C}^n.
\]
\item
Suppose~\(A\) is a Hermitian positive definite matrix split into \(A=C+C^*+D\) where~\(D\) is also Hermitian positive definite. Prove that \(B=C+\omega^{-1}D\) is invertible whenever \(0<\omega<2\). Consider the iteration \(x_{n+1}=x_n+B^{-1}(b-Ax_n)\), with any initial iterate \(x_0\). Prove that \(x_n\) converges to \(x=A^{-1}b\) whenever \(0<\omega<2\).
\item
Let \(f:R\to R\) be a contraction with constant \(\lambda\in(0,1)\). Prove that there exists a unique fixed point~\(\alpha\) of~\(f\). Prove the inequality 
\[
\lvert f^n(x)-\alpha\rvert\leq \frac{\lambda^n}{1-\lambda}\lvert f(x)-x\rvert,\qquad \forall x\in R,\ \forall n\in N.
\]
\item
Derive the two point Gaussian quadrature for approximating \(\int_{-1}^1 f(x)x^2\,dx\).
\item
State the Hermite interpolation problem. Prove that the problem is well-posed, i.e. prove the existence and uniqueness of solutions. Derive the error formula for the interpolating polynomial.
\item
Describe Simpson’s rule for numerical integration (\(n=3\)). State and prove the error formula of Simpson’s rule.
\item
Let \(w(x)>0\) be integrable on \([a,b]\). Define an orthogonal polynomial family \(\{\phi_n\}\) on \([a,b]\) with weight \(w(x)\). Prove that
\begin{examsubparts}
\item[\textnormal{(i)}]
\(\phi_n(x)\) has~\(n\) distinct roots in the interval \((a,b)\), and
\item[\textnormal{(ii)}]
between any two consecutive zeros of \(\phi_n\) there is a zero of \(\phi_{n-1}\).
\end{examsubparts}
\end{examproblems}
\end{document}
